\section{Detection content}
\label{sec:detection}

Eleven rules cover the five scenarios. They fall into three kinds, and the
distinction is not cosmetic --- the kinds behave differently under both attack and
noise.

\begin{description}
\item[Policy rules] fire on a violation of a control that nominal operations
never produce: a command refused by an allowlist, a frame refused by the
vehicle, an audit trail stopped. They are the cheapest and most precise rules in
the pack and, in our results, the ones that carry most of the detection load.
\item[Thresholded rules] aggregate an individually noisy signal over a window: a
burst of failed authentications, a cluster of frame integrity failures, a count
of distinct objects read. Their threshold trades latency against noise directly.
\item[Behavioural rules] compare against a profile learned during a warm-up
window: an unprofiled network origin, a principal issuing its first critical
command, an hourly read volume far above baseline. They are the only rules that
can catch an adversary who violates no policy.
\end{description}

\begin{table}[t]
\centering
\small
\caption{Detection content. Tier determines latency: \texttt{stream} rules see an
event within seconds, \texttt{batch} rules cannot beat their own aggregation
window.}
\label{tab:rules-catalogue}
\begin{tabular}{@{}llll@{}}
\toprule
ID & Kind & Tier & Detects \\
\midrule
R01 & thresholded  & stream & A1 \\
R02 & behavioural  & stream & A1, A2, A3 \\
R03 & policy       & stream & A2 \\
R04 & behavioural  & stream & A2 \\
R05 & policy       & stream & A2, A4 \\
R06 & thresholded  & stream & A3 \\
R07 & behavioural  & batch  & A3 \\
R08 & policy       & stream & A4 \\
R09 & thresholded  & batch  & A5 \\
R10 & statistical  & batch  & A5 \\
R11 & policy       & stream & A3, A5 \\
\bottomrule
\end{tabular}
\end{table}

Every rule maps to a SPARTA technique \cite{sparta} and, where an analogue
exists, to ATT\&CK \cite{attack}. No attack depends on a single rule; that is a
design requirement, and the ablation in \S\ref{sec:results} exists partly to
check that removing one control does not open a blind spot.

\subsection{Two rules worth describing in detail}

\paragraph{R06, frame integrity, and why it is thresholded.}
A radio link corrupts frames on its own, particularly at low elevation at the
edges of a pass. In our model the frame error probability rises exponentially as
elevation falls, producing roughly six to nine corrupted frames per week out of
about \num{1800}. A rule that alerts on a \emph{single} failed tag detects
tampering approximately \SI{60}{\second} sooner and generates enough noise to be
switched off within a week --- the exact failure mode a detection architecture
exists to avoid. R06 therefore requires three failures within
\SI{120}{\second}. \S\ref{sec:results} quantifies both sides of that trade.

\paragraph{R07, physics, and the rail exclusion.}
R07 checks three things about each telemetry channel: that it lies inside its
physical envelope, that it is not changing faster than the subsystem can change,
and that it is not frozen. The third check is what catches a masking attack that
holds a degrading measurement constant. Its first implementation was the largest
single source of false positives in the testbed, because a fully charged battery
in sunlight is legitimately constant. Excluding channels resting against a
physical rail removed that entire class. The lesson generalises beyond this
testbed: a stuck-value detector on any physical process must know which constants
are physics and which are faults.

\paragraph{Analytics should not run on bits already known to be corrupt.}
R07's second-largest false-positive source was subtler than the rail problem and
more instructive. A corrupted frame contains implausible values by definition, so
the physics rule dutifully alerted on link noise that the integrity check had
\emph{already} rejected --- two rules alarming on one event, one of them
uselessly. Excluding samples whose frame failed its tag check removed every such
alert and left a clean division of labour: R06 owns provenance failures, R07 owns
content. The general form of the lesson is that a detection pipeline should
establish trust in an input before reasoning about its meaning.

\subsection{Portability}

Each rule exists in two synchronised forms: executable Python, which both the
experiments and the deployed function import, and Sigma \cite{sigma}, which can
be converted for another backend. A test asserts that the two sets of rule
identifiers are equal, so the packs cannot silently drift apart --- a failure mode
common enough in practice to be worth automating against.
